\section{Orbital transfer and phasing}
\label{sec:transfer}

\subsection{Hohmann geometry}

The minimum-energy two-impulse transfer between coplanar circular orbits is an
ellipse tangent to both \cite{curtis}. With $R_1$ and $R_2$ the two radii,

\begin{equation}
a = \frac{R_1 + R_2}{2}, \qquad
e = \frac{R_2 - R_1}{R_2 + R_1},
\label{eq:ellipse}
\end{equation}

giving $a = \mSMA\unit{km}$ and $e = \mEcc$. Circular and elliptical speeds
follow from the vis-viva equation,

\begin{equation}
v_c(r) = \sqrt{\frac{\mu_M}{r}}, \qquad
v(r) = \sqrt{\mu_M\!\left(\frac{2}{r} - \frac{1}{a}\right)},
\label{eq:visviva}
\end{equation}

and because the transfer is tangent at both ends the two impulses are purely
along-track:

\begin{equation}
\dv_1 = v(R_1) - v_c(R_1), \qquad
\dv_2 = v_c(R_2) - v(R_2).
\label{eq:dv}
\end{equation}

The time of flight is half the ellipse period,
$\Delta t = \pi\sqrt{a^3/\mu_M} = \mTOF\unit{s}$, i.e.\ $\mTOFmin$ minutes.
Table~\ref{tab:hohmann} collects the design.

\begin{table}[htbp]
\centering
\caption{Hohmann transfer design, $100 \to 400\unit{km}$ circular lunar orbit.
The two impulses together cost $\mDVtot\unit{m/s}$, which is the reference
against which every later correction is measured.}
\label{tab:hohmann}
\small
\begin{tabular}{lrl}
\toprule
Quantity & Value & Unit \\
\midrule
Semi-major axis $a$              & \mSMA   & km \\
Eccentricity $e$                 & \mEcc   & -- \\
Circular speed at $R_1$          & 1.6335  & km/s \\
Circular speed at $R_2$          & 1.5145  & km/s \\
Ellipse speed at periapsis       & 1.6940  & km/s \\
Ellipse speed at apoapsis        & 1.4563  & km/s \\
\midrule
Departure impulse $\dv_1$        & \mDVone & m/s \\
Arrival impulse $\dv_2$          & \mDVtwo & m/s \\
\textbf{Total} $\dv$             & \textbf{\mDVtot} & \textbf{m/s} \\
Time of flight                   & \mTOF   & s \\
\bottomrule
\end{tabular}
\end{table}

\subsection{Phasing: why the vehicle waits nine hours}

Designing the ellipse is the easy half. The chaser may only ignite when the
geometry places the target at the arrival point at the arrival time.

During the transfer the chaser sweeps exactly $\pi$ radians while the target
advances $n_2 \Delta t$, with $n_i = \sqrt{\mu_M/R_i^3}$. Rendezvous therefore
requires the target to lead the chaser at ignition by

\begin{equation}
\Delta\theta_H = \pi - n_2 \Delta t = \mLeadDeg^\circ .
\label{eq:leadangle}
\end{equation}

At the epoch the target leads by only $\phi_0 = 10^\circ$. Because the inner
orbit is the faster one, the target's apparent lead \emph{decreases} at the
relative rate $|n_2 - n_1|$: the required $\mLeadDeg^\circ$ will not grow back
from $10^\circ$, it has just been missed. The chaser must wait for the lead to
run down through zero and all the way round:

\begin{equation}
\Delta\theta_{\text{req}} = 2\pi - (\Delta\theta_H - \phi_0), \qquad
t_{\text{wait}} = \frac{\Delta\theta_{\text{req}}}{|n_2 - n_1|} = \mTwait\unit{s}
= \mTwaitH\unit{h}.
\label{eq:twait}
\end{equation}

The implementation carries the general conditional, so a geometry in which the
window has \emph{not} been missed takes the short branch instead. Total mission
duration is $\mTmissionH\unit{h}$, of which the transfer itself is barely an
hour. That ratio is the operationally interesting result of this section:
in a rendezvous of this kind, propellant is cheap and patience is the
constraint.

The rendezvous closes exactly by construction. At arrival the chaser is at
$n_1 t_{\text{wait}} + \pi$ and the target at
$\phi_0 + n_2(t_{\text{wait}} + \Delta t)$; substituting~\eqref{eq:twait} makes
the difference exactly $2\pi$.

\begin{figure}[htbp]
\centering
\includegraphics[width=\textwidth]{fig02_hohmann_geometry.png}
\caption{Transfer geometry at the epoch (left) and at ignition (right). The
target leads by $10^\circ$ at $t=0$ but the manoeuvre needs $\mLeadDeg^\circ$,
so the chaser coasts $\mTwaitH$ hours on the inner orbit before the burn. The
dashed ellipse is the $\mDVtot\unit{m/s}$ transfer.}
\label{fig:geometry}
\end{figure}

\subsection{Numerical propagation}

The full timeline is integrated in Cartesian coordinates with the unperturbed
two-body field,

\begin{equation}
\dot{\mathbf r} = \mathbf v, \qquad
\dot{\mathbf v} = -\mu_M \frac{\mathbf r}{\lVert \mathbf r \rVert^3},
\label{eq:twobody}
\end{equation}

using \texttt{ode45} at relative and absolute tolerances of $10^{-12}$. At
$t = t_{\text{wait}}$ the departure impulse is applied along the instantaneous
velocity unit vector; at apoapsis the second impulse circularises the orbit so
that the constants of motion display a third plateau.

\subsection{Verification against the closed-form solution}
\label{sec:transfer:verif}

The same timeline is then evaluated analytically and differenced. Circular legs
are trivial. On the transfer ellipse the mean anomaly $M = n_H (t - t_{\text{wait}})$
is inverted for the eccentric anomaly by Newton--Raphson,

\begin{equation}
E_{k+1} = E_k - \frac{E_k - e\sin E_k - M}{1 - e\cos E_k},
\label{eq:kepler}
\end{equation}

iterated to $10^{-14}$, and the perifocal state

\begin{equation}
\mathbf r_{pf} = \begin{bmatrix} a(\cos E - e) \\ a\sqrt{1-e^2}\,\sin E \\ 0 \end{bmatrix}
\label{eq:perifocal}
\end{equation}

is rotated into MCI by $R_z(\theta_p)$, with $\theta_p$ the ignition longitude.

Over the whole $\mTmissionH\unit{h}$ timeline the two solutions differ by at most
$\mKeplerErrMm\unit{mm}$ in position, and the numerically propagated rendezvous
closes to $\mMissKep\unit{m}$. This is the single most important number in the
report: it is the licence to believe every perturbed result that follows.

\paragraph{Two different quantities, two different names.}
The millimetre-class figure above is an \emph{integrator residual}: the largest
difference between the \texttt{ode45} state and the closed-form Kepler state at
the same epochs, stored as \texttt{miss\_num\_vs\_kepler}. It is not the
vehicle-to-vehicle separation quoted for the unperturbed control case in
Table~\ref{tab:miss}, which is $\mMissTwoBody\unit{m}$ and is stored as
\texttt{miss\_LM\_MS\_twobody}. That one is an LM--MS distance produced by the
Section~\ref{sec:perturbations} Cowell pipeline with the perturbations switched
off, so it carries that pipeline's working tolerance and both impulse
applications, not the propagator alone. Both are small, they are allowed to
differ by two orders of magnitude, and quoting one for the other would be
wrong. Earlier drafts of this report did exactly that.
When Section~\ref{sec:perturbations} reports a $\mMissJtwo\unit{km}$ miss under
$J_2$, we know that $15\unit{km}$ is physics and not integrator drift, because
the same integrator reproduces the analytical solution to a fifth of a
millimetre.

\begin{figure}[htbp]
\centering
\includegraphics[width=0.93\textwidth]{fig04_position_errors.png}
\caption{Difference between the \texttt{ode45} solution and the closed-form
Kepler solution over the full timeline, log scale. The error grows smoothly to
$\mKeplerErrMm\unit{mm}$, three orders of magnitude below the $1\unit{mm}$
acceptance line. Dashed verticals mark the two impulses.}
\label{fig:keplererr}
\end{figure}

\subsection{Constants of motion}

For a further independent check the specific energy, angular momentum and
eccentricity vector are evaluated along the trajectory,

\begin{equation}
\varepsilon = \frac{v^2}{2} - \frac{\mu_M}{r}, \qquad
\mathbf h = \mathbf r \times \mathbf v, \qquad
\mathbf e = \frac{1}{\mu_M}\!\left[\left(v^2 - \frac{\mu_M}{r}\right)\mathbf r - (\mathbf r \cdot \mathbf v)\mathbf v\right].
\label{eq:constants}
\end{equation}

These must be piecewise constant, stepping only at the two impulses.
Table~\ref{tab:constants} compares the simulated plateaus with the closed-form
values $\varepsilon = -\mu_M/2a$ and $\lVert\mathbf h\rVert = \sqrt{\mu_M a(1-e^2)}$;
they agree to twelve significant figures.

\begin{table}[htbp]
\centering
\caption{Constants of motion per leg. Analytical and simulated values agree to
better than $10^{-12}$ relative, confirming that the impulses are applied
cleanly and that no energy leaks through the integrator.}
\label{tab:constants}
\small
\begin{tabular}{lrrr}
\toprule
Leg & $e$ & $\lVert\mathbf h\rVert$ [km$^2$/s] & $\varepsilon$ [km$^2$/s$^2$] \\
\midrule
Inner circular   & $0$      & $3001.40$ & $-1.33417$ \\
Transfer ellipse & $\mEcc$  & $3112.61$ & $-1.23347$ \\
Outer circular   & $0$      & $3237.17$ & $-1.14691$ \\
\bottomrule
\end{tabular}
\end{table}
